% Generated draft: requires mathematical and source review.
\documentclass[11pt]{article}
\usepackage[margin=25mm]{geometry}
\usepackage{fontspec}
\setmainfont{Latin Modern Roman}
\usepackage{amsmath,amssymb,mathtools,unicode-math}
\setmathfont{Latin Modern Math}
\usepackage{enumitem,needspace,xurl,hyperref}
\setlength{\parindent}{0pt}
\setlength{\parskip}{5pt}
\setlength{\emergencystretch}{3em}
\newcommand{\N}{\mathbb{N}}
\newcommand{\Z}{\mathbb{Z}}
\newcommand{\Q}{\mathbb{Q}}
\newcommand{\R}{\mathbb{R}}
\newcommand{\C}{\mathbb{C}}
\newcommand{\F}{\mathbb{F}}
\newcommand{\A}{\mathbb{A}}
\newcommand{\PP}{\mathbb{P}}
\newcommand{\E}{\mathbb{E}}
\newcommand{\eps}{\varepsilon}
\newcommand{\dd}{\,\mathrm{d}}
\begin{document}
\section*{UnsolvedMath ID: 800860}
\textbf{Source problem code:} catalogue entry
\subsection{Definitions and mathematical statement}
Let \(A=(a_{ij})\) be a matrix whose entries are independent and identically distributed standard Gaussian random variables, as stipulated by the supplied excerpt. Under the usual real Gaussian convention, each entry has distribution \(\mathcal N(0,1)\); the excerpt does not explicitly specify the scalar field, matrix dimensions, or any dimension-dependent normalization.

The source asks: determine \(\omega(A)\). Does this quantity satisfy the bound in the conjecture referred to by the source? If so, what is the parameter \(\beta\)?

The supplied record does not define \(\omega\), reproduce the conjecture or its bound, or define \(\beta\). It also does not specify whether determining the value means finding a distribution, an expectation, an almost-sure value, a high-probability estimate, or an asymptotic quantity. Consequently, these questions cannot be converted into a fully specified mathematical statement from the supplied excerpt alone. The referenced bound and its quantifiers must remain unresolved.

The excerpt's speaker explicitly states that they do not know the state of the art and that a quick search did not reveal an already known answer. This is a supplied qualification, not a verified assertion of openness.
\subsection{Short English statement}
The source asks for \(\omega(A)\) when \(A\) has iid standard Gaussian entries, whether a previously stated conjectural bound holds, and the corresponding \(\beta\). The definitions of \(\omega\) and \(\beta\), the conjecture, and the matrix regime are absent, so a precise formulation requires additional source context.
\subsection{Sources}
\begin{itemize}
\item Local UnsolvedMath registry, permanent ID \(800860\), exact source problem code \texttt{catalogue entry}, supplied title \emph{Randomstrasse101 Open Problems (2025) --- Problem 3}. Upstream revision: \texttt{d371ee917ac9647fbb06cd6fc529deec6f3dda82}.
\item Supplied attribution: Bandeira et al., \emph{Randomstrasse101: Open Problems of 2025} (2026), Problem 3. Supplied source URL: \url{https://arxiv.org/abs/2603.29571}. The record reports access on 2026-10-16 and extraction by \texttt{arxiv\_latex}; neither the access nor the source contents have been independently verified here.
\item The registry supplies the status label \texttt{open} and proposed year 2026. The excerpt itself qualifies its question by reporting uncertainty about the state of the art and a quick search that did not reveal an existing answer. These are supplied source and curation claims, not an independent literature assessment.
\end{itemize}
\end{document}
